% Generated draft: requires mathematical and source review.
\documentclass[11pt]{article}
\usepackage[margin=25mm]{geometry}
\usepackage{fontspec}
\setmainfont{Latin Modern Roman}
\usepackage{amsmath,amssymb,mathtools,unicode-math}
\setmathfont{Latin Modern Math}
\usepackage{enumitem,needspace,xurl,hyperref}
\setlength{\parindent}{0pt}
\setlength{\parskip}{5pt}
\setlength{\emergencystretch}{3em}
\newcommand{\N}{\mathbb{N}}
\newcommand{\Z}{\mathbb{Z}}
\newcommand{\Q}{\mathbb{Q}}
\newcommand{\R}{\mathbb{R}}
\newcommand{\C}{\mathbb{C}}
\newcommand{\F}{\mathbb{F}}
\newcommand{\A}{\mathbb{A}}
\newcommand{\PP}{\mathbb{P}}
\newcommand{\E}{\mathbb{E}}
\newcommand{\eps}{\varepsilon}
\newcommand{\dd}{\,\mathrm{d}}
\begin{document}
\section*{UnsolvedMath ID: 8800005}
\textbf{Source problem code:} catalogue entry
\subsection{Definitions and mathematical statement}
For a finite cyclic group \(G=\langle g\rangle\), the discrete-logarithm problem is to recover \(x\) modulo \(|G|\) from \(g\) and \(h=g^x\). In a finite-field system \(G\) is typically a specified subgroup of a finite field's multiplicative group. The source asks whether Weil descent compromises a proposed realization of such systems in a class group.

The realization itself is not included: no algebraic object defining the class group, map relating the finite-field subgroup to it, parameter family, subgroup order, or public encoding is supplied. These are essential to a statement comparing discrete-logarithm difficulty. Denote the missing construction by \(\mathcal R\); the target is to determine whether a Weil-descent method yields a relevant algorithmic advantage against the discrete-logarithm instances produced by \(\mathcal R\).

For context, Weil restriction for a finite field extension \(K/k\) can represent \(K\)-rational points of a suitable algebraic object by \(k\)-rational points of an object of larger dimension. This geometric reformulation alone does not specify an attack complexity. A precise security question must fix field characteristic and extension degree, the construction and descent, instance distribution, complexity model, and benchmark algorithms; it must also distinguish rigorous bounds, heuristic running times, and concrete parameter security. No universal insecurity statement for class groups or finite-field systems is substituted.

The supplied summary reports Weil descent as a documented attack avenue, especially relevant to some small-characteristic or composite-degree constructions, while saying the specific proposed realization was not conclusively assessed. No exact construction or attack reference is supplied. The 2006 workshop problem and later partial-progress assessment are source attributions pending current primary-literature review.
\subsection{Short English statement}
Does Weil descent make discrete logarithms easier in the workshop's proposed class-group realization of a finite-field system? The realization, parameters and complexity criterion are absent, so the security comparison cannot be made precise from this excerpt alone.
\subsection{Sources}
\begin{itemize}
\item Local registry: \url{data/unsolvedmath/problems.json}, permanent UnsolvedMath ID \(8800005\), source code catalogue entry; supplied statement, background, and research summary, updated 2026-08-06. Attributed to Gerhard Frey.
\item IPAM, \emph{Workshop on Number Theory and Cryptography---Open Problems} (2006), attachment 6219, PDF page 18, Problem 4: \url{https://www.ipam.ucla.edu/programs/workshops/workshop-i-number-theory-and-cryptography-open-problems/?tab=schedule}. Supplied PDF/legacy-PPT visual audit accessed 2026-07-29.
\item The provenance reports inspecting 19 surviving attachments and excluding repeated animation/build material and solved definitions. No separate realization specification or Weil-descent complexity publication is supplied in this record.
\end{itemize}
\end{document}
